%!TeX root=../gardentop.tex
\chapter{The Curtain} 
	
\begin{figure}[t!]
\centering
\includegraphics[width=\linewidth]{july}
\end{figure}

 \lettrine[lines=6]{A}{nd} the secret garden bloomed and bloomed and every morning revealed new miracles. In the robin's nest there were Eggs and the robin's mate sat upon them keeping them warm with her feathery little breast and careful wings. At first she was very nervous and the robin himself was indignantly watchful. Even Dickon did not go near the close-grown corner in those days, but waited until by the quiet working of some mysterious spell he seemed to have conveyed to the soul of the little pair that in the garden there was nothing which was not quite like themselves—nothing which did not understand the wonderfulness of what was happening to them—the immense, tender, terrible, heart-breaking beauty and solemnity of Eggs. If there had been one person in that garden who had not known through all his or her innermost being that if an Egg were taken away or hurt the whole world would whirl round and crash through space and come to an end—if there had been even one who did not feel it and act accordingly there could have been no happiness even in that golden springtime air. But they all knew it and felt it and the robin and his mate knew they knew it.

At first the robin watched Mary and Colin with sharp anxiety. For some mysterious reason he knew he need not watch Dickon. The first moment he set his dew-bright black eye on Dickon he knew he was not a stranger but a sort of robin without beak or feathers. He could speak robin (which is a quite distinct language not to be mistaken for any other). To speak robin to a robin is like speaking French to a Frenchman. Dickon always spoke it to the robin himself, so the queer gibberish he used when he spoke to humans did not matter in the least. The robin thought he spoke this gibberish to them because they were not intelligent enough to understand feathered speech. His movements also were robin. They never startled one by being sudden enough to seem dangerous or threatening. Any robin could understand Dickon, so his presence was not even disturbing.

But at the outset it seemed necessary to be on guard against the other two. In the first place the boy creature did not come into the garden on his legs. He was pushed in on a thing with wheels and the skins of wild animals were thrown over him. That in itself was doubtful. Then when he began to stand up and move about he did it in a queer unaccustomed way and the others seemed to have to help him. The robin used to secrete himself in a bush and watch this anxiously, his head tilted first on one side and then on the other. He thought that the slow movements might mean that he was preparing to pounce, as cats do. When cats are preparing to pounce they creep over the ground very slowly. The robin talked this over with his mate a great deal for a few days but after that he decided not to speak of the subject because her terror was so great that he was afraid it might be injurious to the Eggs.

When the boy began to walk by himself and even to move more quickly it was an immense relief. But for a long time—or it seemed a long time to the robin—he was a source of some anxiety. He did not act as the other humans did. He seemed very fond of walking but he had a way of sitting or lying down for a while and then getting up in a disconcerting manner to begin again.

One day the robin remembered that when he himself had been made to learn to fly by his parents he had done much the same sort of thing. He had taken short flights of a few yards and then had been obliged to rest. So it occurred to him that this boy was learning to fly—or rather to walk. He mentioned this to his mate and when he told her that the Eggs would probably conduct themselves in the same way after they were fledged she was quite comforted and even became eagerly interested and derived great pleasure from watching the boy over the edge of her nest—though she always thought that the Eggs would be much cleverer and learn more quickly. But then she said indulgently that humans were always more clumsy and slow than Eggs and most of them never seemed really to learn to fly at all. You never met them in the air or on tree-tops.

After a while the boy began to move about as the others did, but all three of the children at times did unusual things. They would stand under the trees and move their arms and legs and heads about in a way which was neither walking nor running nor sitting down. They went through these movements at intervals every day and the robin was never able to explain to his mate what they were doing or trying to do. He could only say that he was sure that the Eggs would never flap about in such a manner; but as the boy who could speak robin so fluently was doing the thing with them, birds could be quite sure that the actions were not of a dangerous nature. Of course neither the robin nor his mate had ever heard of the champion wrestler, Bob Haworth, and his exercises for making the muscles stand out like lumps. Robins are not like human beings; their muscles are always exercised from the first and so they develop themselves in a natural manner. If you have to fly about to find every meal you eat, your muscles do not become atrophied (atrophied means wasted away through want of use).

When the boy was walking and running about and digging and weeding like the others, the nest in the corner was brooded over by a great peace and content. Fears for the Eggs became things of the past. Knowing that your Eggs were as safe as if they were locked in a bank vault and the fact that you could watch so many curious things going on made setting a most entertaining occupation. On wet days the Eggs' mother sometimes felt even a little dull because the children did not come into the garden.

But even on wet days it could not be said that Mary and Colin were dull. One morning when the rain streamed down unceasingly and Colin was beginning to feel a little restive, as he was obliged to remain on his sofa because it was not safe to get up and walk about, Mary had an inspiration.

»Now that I am a real boy,« Colin had said, »my legs and arms and all my body are so full of Magic that I can't keep them still. They want to be doing things all the time. Do you know that when I waken in the morning, Mary, when it's quite early and the birds are just shouting outside and everything seems just shouting for joy—even the trees and things we can't really hear—I feel as if I must jump out of bed and shout myself. And if I did it, just think what would happen!«

Mary giggled inordinately.

»The nurse would come running and Mrs~Medlock would come running and they would be sure you had gone crazy and they'd send for the doctor,« she said.

Colin giggled himself. He could see how they would all look—how horrified by his outbreak and how amazed to see him standing upright.

»I wish my father would come home,« he said. »I want to tell him myself. I'm always thinking about it—but we couldn't go on like this much longer. I can't stand lying still and pretending, and besides I look too different. I wish it wasn't raining to-day.«

It was then Mistress Mary had her inspiration.

»Colin,« she began mysteriously, »do you know how many rooms there are in this house?«

»About a thousand, I suppose,« he answered.

»There's about a hundred no one ever goes into,« said Mary. »And one rainy day I went and looked into ever so many of them. No one ever knew, though Mrs~Medlock nearly found me out. I lost my way when I was coming back and I stopped at the end of your corridor. That was the second time I heard you crying.«

Colin started up on his sofa.

»A hundred rooms no one goes into,« he said. »It sounds almost like a secret garden. Suppose we go and look at them. You could wheel me in my chair and nobody would know where we went.«

»That's what I was thinking,« said Mary. »No one would dare to follow us. There are galleries where you could run. We could do our exercises. There is a little Indian room where there is a cabinet full of ivory elephants. There are all sorts of rooms.«

»Ring the bell,« said Colin.

When the nurse came in he gave his orders.

»I want my chair,« he said. »Miss Mary and I are going to look at the part of the house which is not used. John can push me as far as the picture-gallery because there are some stairs. Then he must go away and leave us alone until I send for him again.«

Rainy days lost their terrors that morning. When the footman had wheeled the chair into the picture-gallery and left the two together in obedience to orders, Colin and Mary looked at each other delighted. As soon as Mary had made sure that John was really on his way back to his own quarters below stairs, Colin got out of his chair.

»I am going to run from one end of the gallery to the other,« he said, »and then I am going to jump and then we will do Bob Haworth's exercises.«

And they did all these things and many others. They looked at the portraits and found the plain little girl dressed in green brocade and holding the parrot on her finger.

»All these,« said Colin, »must be my relations. They lived a long time ago. That parrot one, I believe, is one of my great, great, great, great aunts. She looks rather like you, Mary—not as you look now but as you looked when you came here. Now you are a great deal fatter and better looking.«

»So are you,« said Mary, and they both laughed.

They went to the Indian room and amused themselves with the ivory elephants. They found the rose-coloured brocade boudoir and the hole in the cushion the mouse had left but the mice had grown up and run away and the hole was empty. They saw more rooms and made more discoveries than Mary had made on her first pilgrimage. They found new corridors and corners and flights of steps and new old pictures they liked and weird old things they did not know the use of. It was a curiously entertaining morning and the feeling of wandering about in the same house with other people but at the same time feeling as if one were miles away from them was a fascinating thing.

»I'm glad we came,« Colin said. »I never knew I lived in such a big queer old place. I like it. We will ramble about every rainy day. We shall always be finding new queer corners and things.«

That morning they had found among other things such good appetites that when they returned to Colin's room it was not possible to send the luncheon away untouched.

When the nurse carried the tray down-stairs she slapped it down on the kitchen dresser so that Mrs~Loomis, the cook, could see the highly polished dishes and plates.

»Look at that!« she said. »This is a house of mystery, and those two children are the greatest mysteries in it.«

»If they keep that up every day,« said the strong young footman John, »there'd be small wonder that he weighs twice as much to-day as he did a month ago. I should have to give up my place in time, for fear of doing my muscles an injury.«

That afternoon Mary noticed that something new had happened in Colin's room. She had noticed it the day before but had said nothing because she thought the change might have been made by chance. She said nothing to-day but she sat and looked fixedly at the picture over the mantel. She could look at it because the curtain had been drawn aside. That was the change she noticed.

»I know what you want me to tell you,« said Colin, after she had stared a few minutes. »I always know when you want me to tell you something. You are wondering why the curtain is drawn back. I am going to keep it like that.«

»Why?« asked Mary.

»Because it doesn't make me angry any more to see her laughing. I wakened when it was bright moonlight two nights ago and felt as if the Magic was filling the room and making everything so splendid that I couldn't lie still. I got up and looked out of the window. The room was quite light and there was a patch of moonlight on the curtain and somehow that made me go and pull the cord. She looked right down at me as if she were laughing because she was glad I was standing there. It made me like to look at her. I want to see her laughing like that all the time. I think she must have been a sort of Magic person perhaps.«

»You are so like her now,« said Mary, »that sometimes I think perhaps you are her ghost made into a boy.«

That idea seemed to impress Colin. He thought it over and then answered her slowly.

»If I were her ghost—my father would be fond of me,« he said.

»Do you want him to be fond of you?« inquired Mary.

»I used to hate it because he was not fond of me. If he grew fond of me I think I should tell him about the Magic. It might make him more cheerful.«